\chap{7}{CONCLUSION AND FUTURE SCOPE}{0pt}
\vspace{-8pt}
\sect{7.1}{CONCLUSION}
The AI-Based Web Application to Analyse Genetic Mutations and Disease Impact (Evo2Serve) represents a step towards making state-of-the-art genomic artificial intelligence usable by students, researchers and educators. By deploying the Evo~2-7B language model as a secured, serverless GPU service and wrapping it in a browser interface, the project addresses a practical limitation of current research: powerful models for variant effect prediction exist, but they live in notebooks that need a data-centre GPU and programming skills. Our platform lets a user choose a human genome assembly, search a gene such as \textit{BRCA1}, read its reference sequence, list the known ClinVar variants, enter a substitution, deletion, insertion or deletion--insertion, and obtain within seconds a ``likely pathogenic'' or ``likely benign'' label with a confidence value, placed next to the expert ClinVar classification of the same variant.

The technological stack, comprising Python, FastAPI, Modal, PyTorch and FlashAttention on the model side, and Next.js, React, TypeScript, Tailwind CSS, shadcn/ui, Zod, jose, bcrypt and MongoDB on the web side, provides a robust and flexible architecture. Modal allowed the 7-billion-parameter model to be hosted on an H100 GPU that scales between zero and three containers and is billed only while a request is processed, with the model weights cached on a persistent volume. The two-tier design, in which the web tier validates the request and checks the session and the GPU tier verifies the token and validates the request again, keeps the expensive resource protected. Public services from UCSC, NCBI and ClinVar supply the sequence, gene and variant data, so the platform needs no genomic database of its own.

The decision rule converts the difference of the mean log-likelihoods of the mutated and the reference window into a label with a threshold chosen by Youden's index on the \textit{BRCA1} saturation genome editing data. On the 500 calibration variants the zero-shot model reached an AUROC of approximately 0.73. The project placed this number in context by computing the same metric for five established predictors on exactly the same rows: Align-GVGD 0.878, CADD 0.791, phyloP 0.777, PolyPhen-2 0.683 and SIFT 0.656. Evo~2 is clearly better than chance, SIFT, PolyPhen-2 and population frequency on this slice, and below the alignment-based predictors. Since the slice covers one tenth of the assayed interval, the figure is a statement about that region; the published AUROC of the same model on the whole assay is 0.9037.

Because a prediction is only as good as the coordinate and the sequence it is computed on, the project also verified the software around the model, and this is its most distinctive contribution. The calibration data reproduce the published composition of the assay, and the mutation functions produce sequences of the expected length for every variant type. The verification also located eight places where the deployed pipeline differs from the conditions of its calibration: the request carries no reference allele, the deployed window is one base longer than the calibration window, the threshold was calibrated on GRCh37 while the service answers for GRCh38, the browser forwards the GRCh38 coordinate whatever assembly is selected, the two validators disagree on 6 of 55 cases and mishandle 7 more, the reference window is re-scored on every request, a null substitution receives the maximum confidence, and the threshold constants have no archived run artefact. Each of these has a specific and inexpensive repair, and carrying the reference allele in the request closes most of the others at once.

In conclusion, Evo2Serve demonstrates that a large genomic language model can be turned into an accessible, secured and cost-efficient web service, and that verifying the surrounding software is as important as the model itself. The platform is a usable interface for reading a model's opinion about a sequence and for learning how computational evidence relates to curated classification. It is not yet a calibrated variant classifier, its labels are computational scores and not clinical interpretations, and it should be used for education and research triage only. Under the ACMG/AMP framework such evidence can support, but never replace, expert interpretation. The work contributes a deployed system, a reproducible reference for the calibration data and five baseline predictors, and a concrete list of repairs that point the way to a calibrated tool.

The verification study has been released as an open-access preprint on Zenodo \cite{shet2026evo2serve}, under DOI 10.5281/zenodo.24568259, and the source code, analysis scripts and their machine-readable outputs are available in the public project repository at \url{https://github.com/itsvishwaradhya/variant-analysis-evo2}.
\sectend{7.1}

\newpage
\sect{7.2}{FUTURE SCOPE}
Evo2Serve has substantial potential for improvement and extension. The most immediate work is to implement the repairs identified by the verification. The request format should carry the reference allele, and the service should compare it with the fetched reference base and reject a mismatch; this single check catches a coordinate in the wrong assembly, an off-by-one coordinate and a null substitution at once. The browser should forward the position of the assembly that the user selected instead of the GRCh38 position of the ClinVar record, and the service should validate the assembly and chromosome against the list that UCSC publishes and refuse values containing separator characters. The web-tier and GPU-tier validators should be generated from one schema, so that they cannot disagree, and a single window rule should be used by the calibration routine and the endpoint. The reference score of a window is a function of the coordinate only and can be cached across requests, which would remove one of the two forward passes for repeated queries.

Support for further assemblies, for example the complete telomere-to-telomere CHM13 assembly of the human genome, and for genomes of other organisms would widen the use of the platform, provided that the calibration is repeated for each of them. A second priority is recalibration. The present threshold was fitted on 500 substitutions in one region of one gene, without a held-out split, on GRCh37 coordinates and under a window that differs from the deployed one. Future work should recalibrate on all 3,893 variants of the assay with a held-out test split, under the deployed window and in the coordinates of the assembly that the service uses, and should report the operating point with an interval. The confidence value should be replaced by a calibrated probability, following the ClinGen approach of converting a score into evidence strength with local posterior probabilities. Separate thresholds should be calibrated for deletions, insertions and deletion--insertions, whose predictions are currently provisional, and for other genes such as \textit{BRCA2}. Because the published results suggest that the strength of Evo~2 differs between coding and non-coding variants, a single threshold for both is unlikely to be optimal, and the service should return the score and a calibrated probability alongside the label and abstain where the calibration does not apply.

A third direction is the model itself. Evo~2-40B, which needs several GPUs, could be evaluated against the 7B model within the same service. The published results show that a supervised classifier trained on Evo~2 embeddings reaches an AUROC of 0.95 on a held-out split of the \textit{BRCA1} data, which suggests that embeddings extracted from the existing pipeline could be combined with a lightweight trained head, or with the alignment-based predictors that outperform the zero-shot score on the calibration slice, in an ensemble. Because language models trained on sequence alone can reflect phylogenetic and ancestry bias, any such extension should be evaluated across ancestries, and genes of particular interest to Indian populations, which are under-represented in current variant databases, should be added with suitable calibration data.

The platform can also be extended with new functionality. Batch analysis of VCF files, export of a report as a PDF document, and a per-position scan that scores every possible substitution in a window and draws the result as a track along the sequence would turn the single-variant tool into an exploratory one. Explanations in Kannada and English, a guided tutorial that walks a student through the ACMG/AMP framework, and a saved history of analyses for signed-in users, kept only with the consent of the user, would strengthen the educational role. A moderated sharing mode would allow a class to analyse the same set of variants and compare their interpretations.

Finally, the operation of the service needs systematic engineering. Latency, throughput and cost per request were not measured in this project, and a benchmark covering cold starts, warm requests and bursts of concurrent requests would show how the configuration of three containers and a 120-second scale-down window behaves under load. Request batching, rate limiting, structured logging, monitoring of failures and a continuous-integration job that re-runs the verification scripts after every change would make regressions visible before they reach users. Clinical use would require a formal validation study, a quality management system and regulatory review, which are outside the scope of this work; until then the platform should remain a research and teaching tool, and every extension should preserve the wording, flags and limits that keep its output honest.
\sectend{7.2}
